\section{Background and Motivation}\label{sec:motivation}
\subsection{Independent Durable Updates and Finite Append Resources}
Independent file updates provide opportunities for parallel persistence.
Different databases and message streams update their own files, and \texttt{fsync} must confirm each file's contents and recovery metadata.
The filesystem must still allocate new blocks for overwrites, update mappings, and return after durable completion.

ZNS requires sequential appends within each zone and separately limits active and open zones.
Closing a partially written zone releases open occupancy but retains active occupancy.
Our device permits sixteen of each, so 32 workers cannot obtain 32 physical append heads by permanently reserving one writable zone each.
Software request concurrency and physical zone occupancy therefore require separate management.
Foreground commits, ordinary asynchronous writeback, and GC must share this budget.

\subsection{Independent Updates Converge on Shared Allocation State}
F2FS stores file contents in DATA and inode information and block addresses in NODE~\cite{r1}.
An overwrite allocates a new DATA location and records its address in the corresponding NODE.
Both block types occupy the main area (MAIN).
The six-log configuration selects curseg by DATA/NODE and hot, warm, or cold class.
Each curseg records a segment number, summary, and next allocation position.
Files in the same class acquire the same curseg mutex to allocate addresses, record summaries, and advance offsets.

ScaleLFS adds independent GC streams, and Z-LFS expands the zones available to logs.
Ordinary file writeback in both systems retains type-based curseg sharing~\cite{r5,r8,r13,r14}.
Adding reclamation streams or available zones addresses a different layer of concurrency from separating ordinary updates' allocation state.

Files that update independently at the application converge on the same curseg in the allocator.
Figure~\ref{fig:motivation} examines this sharing through throughput, allocation streams, and synchronization paths.
All numerical panels use 64\,KiB updates, with separate runs for throughput, allocation counts, and path diagnostics.
Section~\ref{sec:cost} defines the diagnostic protocols.



As concurrency increases, completion rates for all three baselines rise and then decline.
F2FS, ScaleLFS, and Z-LFS peak at four threads in Figure~\ref{fig:motivation}(a), at 5.82, 6.02, and 3.66\,kIOPS, respectively.
At 32 threads, their rates fall to 0.93, 0.93, and 0.75\,kIOPS.
F2FS allocation counts locate the sharing across files.
At one, two, and four threads, all DATA enters hot DATA and all NODEs enter warm NODE, while the other four streams receive zero blocks (Figure~\ref{fig:motivation}(b)).
Different files thus access the same DATA or NODE allocation state (Figure~\ref{fig:motivation}(d)).

This sharing also appears within complete synchronization calls.
DATA/NODE curseg acquisition and hold intervals occupy 13.3\% of sampled \texttt{fsync} time at one thread, rising to 44.1\%, 72.9\%, and 93.9\% at eight, sixteen, and 32 threads.
At 32 threads, the complete call takes 40.845\,ms, including 38.369\,ms across the two curseg classes.
Acquisition intervals include execution, queuing, and scheduling; Other covers remaining synchronization work and \texttt{balance}.
These measurements locate elapsed time on the shared allocation path.

\subsection{Jointly Organizing Parallel Allocation and Recoverable Commits}\label{sec:joint}
The observations expose different organizational units at the application and allocator.
Applications update independent files, while allocation shares state across a few log classes.
Separating append state by concurrent file update can reduce contention at those entry points.
ZNS limits and file synchronization semantics further determine how these states can operate.

Software append lanes must reuse finite physical resources.
Each commit or asynchronous writeback batch acquires and releases append permission so subsequent updates can reuse the lane.
Maintenance uses a reserved lane under the same physical zone limits.
Partitioning software state by core or thread still requires common accounting for active/open occupancy.

Mapping obligations can also outlive an append lease.
Suppose asynchronous writeback overwrites block $a$ of an existing file at $a'$, clears the DATA dirty flag, and releases the old lease.
The application then overwrites block $b$ and calls \texttt{fsync}.
This synchronization must confirm mappings for both $a'$ and the newly written $b'$.
Inspecting current dirty DATA pages misses $a'$; adding DATA logs alone leaves the corresponding NODEs sharing an allocation entry point.
After distributing NODEs, the original fsync-node sequence no longer directly identifies this commit's terminal.
Physical scan order across lanes also cannot determine commit order.
The commit protocol must connect mapping obligations across leases to selected NODEs and the terminal that confirms them.

Z-ParaPack uses PACK to define each update's commit closure and lanes to separate append state across updates.
Leases manage append permission within the common zone budget.
Terminals and commit numbers identify complete PACKs and their recovery order, while a reserved maintenance lane carries writeback and reclamation outside commits.
